% 自测题2填空2.2.1—2.2.8，印刷51／扫描PDF59页。
\subsection{自测题2：碰撞、旋转与势能填空}

\sourcetest{2.2.1}{51}{59}
\problemtext 质量相同的A、B两物块竖直连接在一轻弹簧两端；
B由支持面C托住处于静止。突然撤去C，求瞬间A与B各自加速度大小。
\workings
撤去前B在弹簧拉力作用下平衡，有弹簧对B向上力$mg$；
对A为向下力$mg$。撤去支持面瞬间弹簧形变量不突变，
A受向下$mg+mg$，B受$mg-mg=0$，故
\[
\boxed{a_A=2g},\qquad\boxed{a_B=0}.
\]
\crosscheck 只讨论\emph{撤去一瞬间}；弹簧伸缩后两物块加速度会变化。

\sourcetest{2.2.2}{51}{59}
\problemtext 小球沿半径$R$的竖直光滑圆环无摩擦滑动；
整个圆环绕其竖直直径以角速率$\omega$转动。
使小球从环底离开并能在偏离底部的平衡位置停留，需要多大角速率？
\workings
以小球偏离底部的角度$\theta$表示位置，
旋转参考系的有效势能
$U_{\rm eff}=mgR(1-\cos\theta)
-\tfrac12m\omega^2R^2\sin^2\theta$。
非底部平衡点满足
$\cos\theta=g/(\omega^2R)$，故需要
\[
\boxed{\omega>\sqrt{\frac gR}}.
\]
\crosscheck 等号处仅为底部临界稳定性；
要求真正\emph{偏离底部}的平衡需严格大于。

\sourcetest{2.2.3}{51}{59}
\problemtext 水平面上两小球完全非弹性碰撞并粘合。
$m_1=2.0$g，$\bm v_1=10\bm i$cm/s，
$m_2=5.0$g，$\bm v_2=(3\bm i+5\bm j)$cm/s，
求共同速度大小及与$x$轴的夹角。
\workings
动量矢量守恒：
\[
\bm v=\frac{2(10\bm i)+5(3\bm i+5\bm j)}{7}
=5\bm i+\frac{25}{7}\bm j\quad(\mathrm{cm/s}).
\]
因此
\[
\boxed{v=\sqrt{25+(25/7)^2}\simeq6.15\,\mathrm{cm/s}},
\qquad
\boxed{\alpha=\arctan(5/7)\simeq35.54^\circ}.
\]
\crosscheck 动能不守恒，须按\emph{矢量}动量分别求$x,y$分量。

\sourcetest{2.2.4}{51}{59}
\problemtext 质量$m$小球以$v_0$追上远大于它质量的挡板，
挡板沿同方向以$v<v_0$运动。垂直弹性碰撞后，
求球的速度大小以及板对球的冲量大小。
\workings
极重挡板速度碰撞前后可视为恒定。
在随板运动的惯性系，球的入射相对速度为$v_0-v$，
弹性反弹后相对速度变成$-(v_0-v)$。
回到地面系，球碰后有向速度为
$v_f=v-(v_0-v)=2v-v_0$，
\[
\boxed{|v_f|=|2v-v_0|},\qquad
\boxed{J=m|v_f-v_0|=2m(v_0-v)}.
\]
\crosscheck 不能直接把碰后有向速度$2v-v_0$
当成非负\emph{速率}；当$v_0>2v$小球反向运动。

\sourcetest{2.2.5}{51}{59}
\problemtext 竖直轻弹簧原长时下端小球$m$恰接触地面，
把弹簧上端缓慢提起，直到球刚要脱离地面。
求外力的功。
\workings
球离地临界时弹簧向上拉力$kx=mg$，
所以伸长$x=mg/k$。
之前球未上升，重力势能不变，
弹簧储能
\[
\boxed{W_{\rm ext}=\frac12kx^2=\frac{m^2g^2}{2k}}.
\]
\crosscheck 这是\emph{准静态}提升，忽略最终动能；
不能把此临界状态与放手自由振动最低点混淆。

\sourcetest{2.2.6}{51}{59}
\problemtext 质点质量$m$在吸引型平方反比中心力
$\bm F=-k\bm e_r/r^2$下以半径$r$作圆周运动，
求速度大小及取无穷远零势能时的机械能。
\workings
向心条件$mv^2/r=k/r^2$，得
\[
\boxed{v=\sqrt{\frac{k}{mr}}}.
\]
势能$U=-k/r$，动能$K=k/(2r)$，
所以
\[
\boxed{E=K+U=-\frac{k}{2r}}.
\]
\crosscheck 圆轨道束缚态总能量为负，源题指定无穷远处$U=0$。

\sourcetest{2.2.7}{51}{59}
\problemtext 质量$m$的地球卫星在地面上方高度$2R$绕地球作圆轨道，
地球半径$R$、质量$M$、万有引力常数$G$，
求卫星动能及引力势能。
\workings
轨道半径是$r=R+2R=3R$，
圆轨道动能$K=GMm/(2r)$、势能$U=-GMm/r$，
故
\[
\boxed{K=\frac{GMm}{6R}},\qquad
\boxed{U=-\frac{GMm}{3R}}.
\]
\crosscheck 不可将高度$2R$误作圆轨道半径。

\sourcetest{2.2.8}{51}{59}
\problemtext 一圆盘以恒角加速度转动，
初角速度$20\pi$rad/s，转过60转后角速度$30\pi$rad/s，
求角加速度与经过时间。
\workings
角位移$\Delta\theta=60(2\pi)=120\pi$rad，
$\omega_f^2-\omega_0^2=2\alpha\Delta\theta$，故
\[
\boxed{\alpha=\frac{(30\pi)^2-(20\pi)^2}{240\pi}
=\frac{25\pi}{12}\,\mathrm{rad/s^2}}.
\]
时间为
\[
\boxed{\Delta t=\frac{\omega_f-\omega_0}{\alpha}
=\frac{10\pi}{25\pi/12}=4.80\,\mathrm s}.
\]
\crosscheck 60转要换成120π弧度，不能直接将60代入角位移公式。
